\chapter{Основы языка Verilog}
\label{ch:verilog}

В главе~\ref{ch:gates} мы собирали схемы из элементов и записывали их булевыми
формулами. Пока схема маленькая, её можно нарисовать мышью, но сотню элементов
так не соберёшь. Нужен способ записать схему \emph{текстом}, чтобы Logra сама
превратила его в элементы и показала результат. Эта глава~--- основы такого
языка, Verilog. Дальше он будет сопровождать каждую тему: сначала схема и
таблица, потом тот же узел кодом, запуск в Logra и сравнение.

\section{Код описывает схему, а не последовательность действий}

Главное, что нужно понять перед изучением Verilog: \textbf{это не язык
программирования в привычном смысле}. Строчки кода на Verilog не выполняются
одна за другой. Каждая строчка описывает кусок схемы, и все куски
работают \emph{одновременно}, как элементы на макетной плате.

\begin{figure}[H]
\centering
\begin{tikzpicture}
  \node[draw=FpgaBlue, fill=CodeBack, thick, rounded corners, align=left, font=\ttfamily\small, inner sep=3mm] (code) at (0,0) {
  \textcolor{CodeKeyword}{\textbf{assign}} x = a \& b;\\
  \textcolor{CodeKeyword}{\textbf{assign}} y = c | d;\\
  \textcolor{CodeKeyword}{\textbf{assign}} z = x \^{} y;};
  \node[lbl, above=1mm of code] {код на Verilog};
  \begin{scope}[xshift=5.6cm, yshift=-0.2cm]
    \node[gAND] (g1) at (0,0.9) {};
    \node[gOR] (g2) at (0,-0.9) {};
    \node[gXOR] (g3) at (3.0,0) {};
    \draw[wire] (g1.input 1) -- ++(-0.6,0) node[left] {$a$};
    \draw[wire] (g1.input 2) -- ++(-0.6,0) node[left] {$b$};
    \draw[wire] (g2.input 1) -- ++(-0.6,0) node[left] {$c$};
    \draw[wire] (g2.input 2) -- ++(-0.6,0) node[left] {$d$};
    \draw[wire] (g1.output) -| ([xshift=-4mm]g3.input 1) -- (g3.input 1);
    \node[lbl, above] at ($(g1.output)+(0.35,0)$) {x};
    \draw[wire] (g2.output) -| ([xshift=-4mm]g3.input 2) -- (g3.input 2);
    \node[lbl, below] at ($(g2.output)+(0.35,0)$) {y};
    \draw[wire] (g3.output) -- ++(0.6,0) node[right] {$z$};
    \node[lbl] at (1.5,1.9) {схема, которую он описывает};
  \end{scope}
  \draw[arr, FpgaOrange] (code.east) -- (4.1,0);
\end{tikzpicture}
\caption{Три строчки кода --- три логических элемента, работающих одновременно}
\label{fig:code_to_schem}
\end{figure}

Если поменять строчки местами, схема не изменится: провод \code{x} всё равно
соединит выход элемента И со входом «Исключающего ИЛИ».

\begin{diybox}[Порядок строк не важен]
\textbf{Задание.} Проверьте на рис.~\ref{fig:code_to_schem}, что строки \code{assign} работают одновременно.\\
\textbf{В Logra.} «Создать проект» $\to$ «С нуля», модуль \code{top}; вкладка «Код»: входы \code{a, b, c, d}, выход \code{z} и три строки \code{assign} (\code{x = a \& b;}, \code{y = c | d;}, \code{z = x \^{} y;}, провода \code{x} и \code{y} объявите как \code{wire}); «Синтезировать».\\
\textbf{Ожидаемый результат.} На вкладке RTL три элемента: И, ИЛИ и XOR.\\
\textbf{Как проверить.} Поменяйте строки местами, пересоберите: схема на RTL должна остаться той же. \stars{2} Уберите \code{wire}: прочитайте сообщение в «Диагностике» и объясните его.
\end{diybox}

\section{Модуль}

Любая схема на Verilog оформляется как \term{модуль} (\code{module}). Модуль
похож на микросхему: у него есть имя, входы, выходы и «начинка».

\begin{figure}[H]
\centering
\begin{minipage}{0.52\textwidth}
\tcbset{codestyle/.append style={unbreakable}}
\begin{vcode}
module half_adder (   // имя модуля
    input  wire a,    // входы
    input  wire b,
    output wire s,    // выходы
    output wire c
);
    assign s = a ^ b; // сумма
    assign c = a & b; // перенос
endmodule
\end{vcode}
\end{minipage}\hfill
\begin{minipage}{0.44\textwidth}
\centering
\begin{tikzpicture}
  \node[draw=FpgaBlue, fill=FpgaLight, very thick, minimum width=26mm, minimum height=24mm, font=\sffamily\small, align=center] (m) {half\_adder};
  \draw[wire] ([yshift=5mm]m.west) -- ++(-0.9,0) node[left] {a};
  \draw[wire] ([yshift=-5mm]m.west) -- ++(-0.9,0) node[left] {b};
  \draw[wire] ([yshift=5mm]m.east) -- ++(0.9,0) node[right] {s};
  \draw[wire] ([yshift=-5mm]m.east) -- ++(0.9,0) node[right] {c};
  \node[lbl, below=1mm of m] {модуль = «микросхема»};
\end{tikzpicture}
\end{minipage}
\caption{Модуль полусумматора и его «корпус»}
\end{figure}

Разберём ключевые слова:
\begin{itemize}
  \item \code{module имя ( ... ); ... endmodule} --- начало и конец описания;
  \item \code{input}, \code{output}, \code{inout} --- направление порта;
  \item \code{wire} --- провод (соединение);
  \item \code{assign} --- «подключить к проводу слева выход схемы справа».
\end{itemize}

\begin{diybox}[Свой полусумматор]
\textbf{Задание.} Опишите модуль полусумматора из раздела и проверьте его.\\
\textbf{В Logra.} Вкладка «Код»: модуль \code{half\_adder} (имя можно оставить \code{top}); «Синтезировать»; вкладки RTL и «Диаграммы» $\to$ «Перебрать входы».\\
\textbf{Ожидаемый результат.} В RTL один XOR и один AND; на диаграмме $s=a\oplus b$, $c=a\cdot b$ для всех четырёх комбинаций.\\
\textbf{Как проверить.} Сверьте диаграмму с таблицей полусумматора из главы~\ref{ch:arith}; щёлкните по XOR на RTL и убедитесь, что «Инспектор» называет строку \code{assign s = a \^{} b;}. \stars{2} Сделайте вход \code{a} двухразрядным и посмотрите, как изменятся порты.
\end{diybox}

\section{Сигналы: wire и reg}

\begin{description}
  \item[wire] --- провод. Не хранит значение, а лишь передаёт то, что на него
    подаёт какая-то схема. Присваивается через \code{assign}.
  \item[reg] --- переменная, которой присваивают значение внутри блока
    \code{always}. Несмотря на название, \code{reg} становится триггером
    \emph{не всегда}: это зависит от того, в каком \code{always} он описан (см.
    раздел~\ref{sec:always}).
\end{description}

\subsection{Шины (многоразрядные сигналы)}

Несколько проводов можно объединить в \term{шину}:
\begin{vcode}
wire [7:0] data;      // 8 проводов: data[7], data[6], ..., data[0]
reg  [5:0] counter;   // 6-разрядный регистр
wire       bit3 = data[3];      // один провод из шины
wire [3:0] low  = data[3:0];    // младшие 4 провода
\end{vcode}

\begin{figure}[H]
\centering
\begin{tikzpicture}
  \foreach \i in {0,...,7} {
    \pgfmathtruncatemacro{\n}{7-\i}
    \node[draw=FpgaBlue, fill=\ifnum\n<4 orange!20\else FpgaLight\fi, minimum width=9mm, minimum height=8mm, font=\ttfamily\small] (b\i) at (\i*0.9,0) {\n};
  }
  \node[left=2mm of b0, font=\ttfamily\small] {data =};
  \draw[decorate, decoration={brace, amplitude=5pt, mirror}, thick] (b4.south west) -- (b7.south east) node[midway, below=6pt, lbl] {data[3:0]};
  \node[above=1mm of b0, lbl] {старший (MSB)};
  \node[above=1mm of b7, lbl] {младший (LSB)};
\end{tikzpicture}
\caption{Шина \code{[7:0]}: нумерация разрядов справа налево}
\end{figure}

\subsection{Массивы: память}

Если после имени поставить диапазон индексов, получится массив слов:
\code{reg [7:0] mem [0:15];}~--- шестнадцать слов по восемь бит, обращение
\code{mem[addr]}. Так описывают память (ПЗУ, ОЗУ, очередь FIFO); читают и
пишут её внутри \code{always}. Пример с массивом~--- очередь в
главе~\ref{ch:fsmdesign}.

\subsection{Запись чисел}

Числа в Verilog записываются в формате \code{<разрядность>'<система><значение>}:

\begin{table}[H]
\centering
\caption{Запись чисел}
\begin{tabular}{l l l}
\toprule
\textbf{Запись} & \textbf{Значение} & \textbf{Пояснение} \\
\midrule
\code{4'b1010}        & 10  & 4 бита, двоичная система \\
\code{8'hFF}          & 255 & 8 бит, шестнадцатеричная \\
\code{6'd42}          & 42  & 6 бит, десятичная \\
\code{1'b0}, \code{1'b1} & 0, 1 & один бит \\
\code{27\_000\_000}   & 27000000 & подчёркивания игнорируются, их ставят для удобства чтения \\
\code{4'bx01z}        & ---  & \code{x} --- неизвестно, \code{z} --- отключено (высокий импеданс) \\
\bottomrule
\end{tabular}
\end{table}

\begin{diybox}[Шины и числа]
\textbf{Задание.} Сделайте учебные упражнения про провода и шины.\\
\textbf{В Logra.} Вкладка «Обучение», курс «Б. Verilog с нуля»: уроки «Первый провод», «Число из четырёх бит», «Верхняя половина шины».\\
\textbf{Ожидаемый результат.} Все проверки каждого урока зелёные.\\
\textbf{Как проверить.} Затем в «Коде» напишите \code{wire [3:0] low = data[3:0];} для 8-разрядной шины \code{data} и убедитесь по RTL, что выход берёт четыре младших провода. \stars{2} Запишите в \code{4'bx01z} другие значения \code{x} и \code{z} и посмотрите, как их показывает диаграмма.
\end{diybox}

\section{Операторы}

\begin{table}[htbp]
\centering
\caption{Основные операторы Verilog}
\label{tab:ops}
\small
\begin{tabularx}{\textwidth}{l l X}
\toprule
\textbf{Группа} & \textbf{Операторы} & \textbf{Пример и смысл} \\
\midrule
Побитовые & \code{\textasciitilde{} \& | \^{} \textasciitilde\&  \textasciitilde|} &
  \code{a \& b} --- И над каждой парой битов; \code{\textasciitilde{}a} --- НЕ \\
Логические & \code{! \&\& ||} & \code{(x > 3) \&\& en} --- результат 1 бит (истина/ложь) \\
Свёртка & \code{\&a  |a  \^{}a} & \code{\&a} --- И всех битов шины (1, если все единицы) \\
Арифметика & \code{+ - * / \%} & \code{cnt + 1'b1} --- сумматор \\
Сравнение & \code{== != < > <= >=} & \code{cnt == 25} --- компаратор \\
Сдвиги & \code{<< >>} & \code{x << 1} --- сдвиг влево на 1 \\
Склейка & \code{\{a, b\}} & \code{\{4'b0, x\}} --- склеить шины в одну \\
Повтор & \code{\{n\{a\}\}} & \code{\{6\{1'b1\}\}} = \code{6'b111111} \\
Условие & \code{c ? a : b} & мультиплексор: если \code{c}, то \code{a}, иначе \code{b} \\
\bottomrule
\end{tabularx}
\end{table}

Каждый оператор --- это кусок аппаратуры. \code{a + b} --- это сумматор,
\code{a == b} --- компаратор, \code{c ? a : b} --- мультиплексор.

\begin{figure}[H]
\centering
\begin{tikzpicture}
  \node[draw=FpgaDark, fill=orange!15, thick, trapezium, trapezium angle=65, shape border rotate=270, minimum width=20mm, minimum height=10mm, font=\sffamily\small] (mux) at (0,0) {MUX};
  \draw[wire] ([yshift=3mm]mux.west) -- ++(-1,0) node[left] {a};
  \node[font=\footnotesize, anchor=west] at ([yshift=3mm, xshift=1mm]mux.west) {1};
  \draw[wire] ([yshift=-3mm]mux.west) -- ++(-1,0) node[left] {b};
  \node[font=\footnotesize, anchor=west] at ([yshift=-3mm, xshift=1mm]mux.west) {0};
  \draw[wire] (mux.east) -- ++(1,0) node[right] {y};
  \draw[wire, FpgaBlue] (mux.south) -- ++(0,-0.6) node[below] {c};
  \node[font=\ttfamily\small, anchor=west] at (3,0) {assign y = c ? a : b;};
\end{tikzpicture}
\caption{Тернарный оператор \code{?:} --- это мультиплексор}
\end{figure}

\begin{diybox}[Оператор как узел схемы]
\textbf{Задание.} Увидьте, что каждому оператору соответствует свой узел.\\
\textbf{В Logra.} «Код»: порты \code{a, b} (4 бита), \code{c}, выходы \code{s = a + b}, \code{eq = (a == b)}, \code{m = c ? a : b}; «Синтезировать»; RTL. Для упражнений откройте «Обучение»: «Не «И»» и «Склей биты».\\
\textbf{Ожидаемый результат.} На RTL сумматор, компаратор и мультиплексор, как на рис. «Тернарный оператор».\\
\textbf{Как проверить.} Подайте на «Диаграммах» значения и проверьте в уме одно сложение и одно сравнение. \stars{2} Замените \code{c ? a : b} на запись через \code{if} в \code{always @(*)} и убедитесь, что схема не изменилась.
\end{diybox}

\section{Блоки always}
\label{sec:always}

Описывать сложную логику одними \code{assign} неудобно. Для этого есть блок
\code{always}. Он бывает двух видов, и важно их не путать.

\subsection{Комбинационный: always @(*)}

Звёздочка означает: «пересчитывай выход при любом изменении любого входа». Такой
блок описывает \textbf{комбинационную схему без памяти}.
\begin{vcode}
reg [1:0] y;
always @(*) begin
    case (sel)
        2'd0: y = a;
        2'd1: y = b;
        2'd2: y = c;
        default: y = 2'b00;
    endcase
end
\end{vcode}

\subsection{Последовательностный: always @(posedge clk)}

Этот блок срабатывает только по \textbf{переднему фронту тактового сигнала}.
Все \code{reg}, которым присваивают значение внутри него, становятся
\textbf{триггерами}.
\begin{vcode}
reg [3:0] cnt = 0;
always @(posedge clk) begin
    cnt <= cnt + 1'b1;   // каждый такт прибавляем единицу
end
\end{vcode}

\begin{figure}[H]
\centering
\begin{tikzpicture}
  \node[draw=FpgaGreen, fill=green!8, thick, minimum width=28mm, minimum height=18mm, font=\sffamily\small, align=center] (reg) at (4.2,0) {регистр\\\footnotesize 4 бита};
  \node[font=\sffamily\footnotesize, anchor=west] at (reg.west) {D};
  \node[font=\sffamily\footnotesize, anchor=east] at (reg.east) {Q};
  \draw[thick] ([xshift=-1.2mm]reg.south) -- ([yshift=1.8mm]reg.south) -- ([xshift=1.2mm]reg.south);
  \draw[wire] (reg.south) -- ++(0,-0.8) node[below] {clk};
  \node[draw=FpgaBlue, fill=FpgaLight, thick, circle, minimum size=11mm, font=\sffamily\bfseries] (add) at (1.2,0) {$+$};
  \draw[bus, -{Stealth}] (add.east) -- node[above, lbl] {cnt+1} (reg.west);
  \draw[bus, -{Stealth}] (reg.east) -- ++(1.6,0) node[right, text=FpgaDark] {cnt};
  \draw[bus, -{Stealth}] ($(reg.east)+(0.8,0)$) -- ++(0,1.5) -| (add.north);
  \fill[FpgaBlue] ($(reg.east)+(0.8,0)$) circle (1.6pt);
  \draw[wire] (add.south) -- ++(0,-0.7) node[below] {1};
\end{tikzpicture}
\caption{Счётчик: сумматор и регистр с обратной связью. Код \code{cnt <= cnt + 1}
описывает именно эту схему}
\label{fig:counter}
\end{figure}

\begin{figure}[H]
\centering
\begin{tikztimingtable}[timing/slope=0.1, xscale=2.2, yscale=1.4, timing/font=\sffamily\small, timing/rowdist=3.2ex]
  clk & 9{1L 1H} \\
  cnt & 1D{0} 2D{1} 2D{2} 2D{3} 2D{4} 2D{5} 2D{6} 2D{7} 2D{8} 1D{9} \\
\end{tikztimingtable}
\caption{Временная диаграмма счётчика: значение увеличивается по каждому
переднему фронту}
\end{figure}

\subsection{Блокирующее (=) и неблокирующее (<=) присваивание}

\begin{warnbox}[Золотое правило]
\begin{itemize}
  \item В \code{always @(posedge clk)} используйте только \code{<=}.
  \item В \code{always @(*)} используйте только \code{=}.
\end{itemize}
\end{warnbox}

Почему? Неблокирующее присваивание \code{<=} означает: «все правые части
вычисляются по значениям \emph{до} фронта, и все триггеры обновляются
одновременно». Именно так ведут себя настоящие триггеры. Посмотрим на
пример, в котором два регистра обмениваются значениями:

\begin{figure}[H]
\centering
\begin{minipage}{0.47\textwidth}
\tcbset{codestyle/.append style={unbreakable}}
\begin{vcode}
always @(posedge clk) begin
    a <= b;
    b <= a;   // обмен работает!
end
\end{vcode}
\centering
\begin{tikzpicture}[ffn/.style={draw=FpgaGreen, fill=green!8, thick, minimum width=10mm, minimum height=10mm, font=\sffamily\small}]
  \node[ffn] (A) at (0,0) {a};
  \node[ffn] (B) at (2.4,0) {b};
  \draw[arr] (A.north) to[bend left=35] (B.north);
  \draw[arr] (B.south) to[bend left=35] (A.south);
\end{tikzpicture}
\end{minipage}\hfill
\begin{minipage}{0.47\textwidth}
\tcbset{codestyle/.append style={unbreakable}}
\begin{vcode}
always @(posedge clk) begin
    a = b;
    b = a;    // ошибка: b = b
end
\end{vcode}
\centering
\begin{tikzpicture}[ffn/.style={draw=FpgaRed, fill=red!5, thick, minimum width=10mm, minimum height=10mm, font=\sffamily\small}]
  \node[ffn] (A) at (0,0) {a};
  \node[ffn] (B) at (2.4,0) {b};
  \draw[arr] (B.north) to[bend right=35] (A.north);
  \draw[arr] (B.east) to[out=0, in=-60, looseness=6] (B.south east);
\end{tikzpicture}
\end{minipage}
\caption{Слева --- правильный обмен (\code{<=}), справа --- ошибка (\code{=}):
к моменту второй строки \code{a} уже равно \code{b}}
\end{figure}

\begin{diybox}[Счётчик и обмен регистров]
\textbf{Задание.} Опишите 4-разрядный счётчик и разберите ловушку блокирующего присваивания.\\
\textbf{В Logra.} «Код»: \code{reg [3:0] cnt = 0; always @(posedge clk) cnt <= cnt + 1'b1;} с портами \code{clk}, \code{reset}; «Синтезировать»; RTL и «Диаграммы». Упражнения: «Первая ячейка памяти», «Две ячейки друг за другом».\\
\textbf{Ожидаемый результат.} На RTL регистр, сумматор и обратная связь (рис.~\ref{fig:counter}); на диаграмме значение растёт на единицу по каждому фронту и после 15 возвращается к 0.\\
\textbf{Как проверить.} \stars{2} Сделайте два регистра \code{a} и \code{b} с обменом \code{a <= b; b <= a;}, затем замените \code{<=} на \code{=} и сравните диаграммы: во втором случае оба регистра сравняются.
\end{diybox}

\textbf{Образец (счётчик с \code{reset} и \code{enable}).}\par\nopagebreak
\begin{vcode}
module cnt4 (
    input            clk, reset, enable,
    output reg [3:0] cnt
);
    always @(posedge clk) begin
        if (reset)       cnt <= 4'd0;         // сброс важнее счёта
        else if (enable) cnt <= cnt + 1'b1;   // иначе cnt хранит значение
    end
endmodule
\end{vcode}


\section{Условия: if и case}

Внутри \code{always} можно использовать условия. Они тоже превращаются в
мультиплексоры.
\begin{vcode}
always @(posedge clk) begin
    if (reset)
        cnt <= 0;               // синхронный сброс
    else if (enable)
        cnt <= cnt + 1'b1;      // счёт
    // иначе cnt сохраняет значение
end
\end{vcode}

\begin{warnbox}[Защёлки (latch)]
Если в \code{always @(*)} при каком-то условии переменной ничего не присвоено,
синтезатор вынужден «запомнить» старое значение и создаёт \term{защёлку} ---
нежелательный элемент памяти без такта. Всегда пишите \code{else} и
\code{default}!
\end{warnbox}

\begin{diybox}[Мультиплексор и защёлка]
\textbf{Задание.} Опишите выбор по \code{case} и намеренно получите защёлку.\\
\textbf{В Logra.} «Код»: \code{always @(*) case (sel) ... endcase} из раздела про \code{always}; сначала с \code{default}, затем без него и без одного значения \code{sel}; «Синтезировать»; RTL. Упражнения: «Выбор по условию», «Выбор из списка».\\
\textbf{Ожидаемый результат.} С \code{default} на RTL мультиплексор; без \code{default} появляется элемент LATCH (в Logra это видно рядом с DFF).\\
\textbf{Как проверить.} Сравните с предупреждением «Защёлки (latch)» и верните \code{default}. \stars{2} Опишите приоритетный шифратор через \code{if / else if} и посмотрите, чем он отличается по RTL от \code{case}.
\end{diybox}

\section{Приоритет и casez: шифратор}
\label{sec:v-prio}

Шифратор $4\to2$ выдаёт номер входа, на котором стоит единица. Если единиц
несколько, нужен \emph{приоритет}: например, выигрывает старший вход. В
Verilog приоритет задаёт \emph{порядок проверок}: ветки \code{if / else if}
проверяются сверху вниз, строки \code{casez} тоже, и сработает первая
подходящая. В \code{casez} символ \code{?} означает «любое значение»: строка
\code{4'b01??} читается как «бит~3 равен нулю, бит~2 равен единице, остальные
безразличны». Ниже одна и та же функция записана тремя способами: через
\code{if}, через \code{casez} и вентилями.

\vfile{\codepath/05_verilog/enc4.v}{code/05\_verilog/enc4.v}

\textbf{Как читается.} В первом модуле значение \code{valid} сначала
выставляется в единицу, а ветка \code{else} в конце обнуляет его, если ни один
вход не активен: так переменные \code{y} и \code{valid} получают значение при
\emph{любом} наборе входов, и защёлки нет. Во втором модуле роль \code{else}
играет \code{default}. Третий модуль не содержит приоритета в явном виде: он
записывает ту же таблицу истинности формулами, и на вкладке RTL его схема
оказывается из вентилей, а не из мультиплексоров.

Проверяет модуль тестбенч \code{tb\_enc4.v}: он перебирает все 16 входных слов
и убеждается, что три записи дают одинаковый результат (разбор тестбенча ---
в разделе~\ref{sec:v-tb}).

\vfile{\codepath/05_verilog/tb_enc4.v}{code/05\_verilog/tb\_enc4.v}

\begin{diybox}[Приоритетный шифратор]
\textbf{Задание.} Построить шифратор тремя способами и сравнить схемы.\\
\textbf{В Logra.} Вкладка «Код»: вставьте \code{enc4.v}, нажмите «Синтезировать»
(Ctrl+Enter); по очереди выберите верхним модулем \code{enc4\_if},
\code{enc4\_casez} и \code{enc4\_gates}; вкладка RTL, затем «Диаграммы»
$\to$ «Перебрать входы». Тот же тестбенч запускается в терминале:
\code{iverilog -g2012 -o tb tb\_enc4.v enc4.v}, затем \code{vvp tb} --- должна
появиться строка \code{enc4: OK}.\\
\textbf{Ожидаемый результат.} На диаграмме при \code{d = 0110} выход \code{y = 2}
(выигрывает старший), при \code{d = 0000} сигнал \code{valid = 0}.\\
\textbf{Как проверить.} Найдите на RTL разницу между вариантами: у \code{if}
и \code{casez} цепочка мультиплексоров, у третьего модуля несколько вентилей.
\stars{2} Допишите модуль так, чтобы приоритет был у \emph{младшего} входа.
\end{diybox}

\textbf{Образец ответа к \stars{2}.}\par\nopagebreak
\begin{vcode}
module enc4_low (
    input      [3:0] d,
    output reg [1:0] y,
    output reg       valid
);
    always @(*) begin
        valid = 1'b1;
        casez (d)
            4'b???1: y = 2'd0;      // младший вход главнее всех
            4'b??10: y = 2'd1;
            4'b?100: y = 2'd2;
            4'b1000: y = 2'd3;
            default: begin y = 2'd0; valid = 1'b0; end
        endcase
    end
endmodule
\end{vcode}


\section{Мультиплексор и семисегментный дешифратор}
\label{sec:v-mux-seg}

Мультиплексор выбирает один из входов по адресу. Его можно записать
вложенным условием \code{?:} или оператором \code{case}. Оба способа дают одну
и ту же схему, но \code{case} легче читать, когда входов много.

\vfile{\codepath/05_verilog/mux4.v}{code/05\_verilog/mux4.v}

В модуле \code{mux4\_case} последняя строка названа \code{default}, а не
\code{2'd3}: так выход определён при любых значениях \code{s}, даже
неизвестных \code{x} в симуляции. Тестбенч \code{tb\_mux4.v} сверяет оба модуля
с индексацией \code{d[s]} для всех 64 сочетаний входов.

Семисегментный индикатор показывает цифру семью сегментами, обозначенными
буквами \code{a}\ldots\code{g}. Дешифратор превращает четырёхбитовое число в
семь управляющих сигналов. Это таблица истинности, записанная оператором
\code{case}: каждой строке соответствует одна цифра.

\vfile{\codepath/05_verilog/seg7.v}{code/05\_verilog/seg7.v}

Число \code{7'b0111111} читается справа налево как сегменты \code{a, b, c, d,
e, f} (горят) и \code{g} (не горит): получается «0». Если индикатор включается
нулём (общий анод), результат инвертируют: \code{assign seg\_n = \textasciitilde seg;}.

\begin{diybox}[Выбор и индикатор]
\textbf{Задание.} Собрать мультиплексор и дешифратор и проверить их таблицы.\\
\textbf{В Logra.} Вкладка «Код»: \code{mux4.v} и \code{seg7.v} (по одному),
«Синтезировать», RTL, «Диаграммы». Тестбенчи \code{tb\_mux4.v} и
\code{tb\_seg7.v} запускаются в терминале так же, как \code{tb\_enc4.v}.\\
\textbf{Ожидаемый результат.} Мультиплексор повторяет вход \code{d[s]}; для
\code{digit = 8} горят все семь сегментов, для \code{digit = 1} только два.\\
\textbf{Как проверить.} «Диаграммы» $\to$ «Перебрать входы» для мультиплексора;
для дешифратора задайте \code{digit} от 0 до 15 и сверьтесь с таблицей в коде.
\stars{2} Допишите дешифратор для двухбитового числа (цифры 0--3; все остальные
значения невозможны, но \code{default} всё равно нужен).
\end{diybox}

\textbf{Образец ответа к \stars{2}.}\par\nopagebreak
\begin{vcode}
module seg4 (
    input      [1:0] d,
    output reg [6:0] seg
);
    always @(*) begin
        case (d)
            2'd0:    seg = 7'b0111111;
            2'd1:    seg = 7'b0000110;
            2'd2:    seg = 7'b1011011;
            default: seg = 7'b1001111;   // цифра 3
        endcase
    end
endmodule
\end{vcode}


\section{Параметры и константы}

\begin{vcode}
localparam integer CLK_HZ = 27_000_000;  // константа внутри модуля
localparam [1:0]   IDLE = 2'd0, RUN = 2'd1;

module blink #(parameter integer HZ = 1) (...);  // параметр можно
                                                 // менять снаружи
\end{vcode}

\begin{diybox}[Параметр частоты]
\textbf{Задание.} Измените схему одним числом, не трогая логики.\\
\textbf{В Logra.} Проект «Бегущий огонь», строка \code{localparam DIV = 27\_000\_000 / 8;}; урок «Настройка в одном месте» в курсе «Обучение».\\
\textbf{Ожидаемый результат.} После замены \code{8} на \code{2} огонь бежит в четыре раза медленнее (две смены в секунду).\\
\textbf{Как проверить.} «Синтезировать», вкладка «Лаборатория», список «Скорость» в панели сверху позволяет сравнить на глаз. \stars{2} Преобразуйте \code{localparam} в параметр модуля \code{\#(parameter DIV = ...)} и проверьте, что схема собирается так же.
\end{diybox}

\textbf{Образец.} Параметр \code{DIV} меняют при создании экземпляра
(\code{blink \#(.DIV(27\_000\_000 / 2)) u (...)}) или в заголовке:
\begin{vcode}
module blink #(parameter integer DIV = 27_000_000 / 8) (
    input      clk,
    output reg led = 1'b0
);
    reg [24:0] cnt = 25'd0;
    always @(posedge clk) begin
        if (cnt == DIV - 1) begin
            cnt <= 25'd0;
            led <= ~led;                // смена раз в DIV тактов
        end else cnt <= cnt + 25'd1;
    end
endmodule
\end{vcode}


\section{Циклы for, generate и функции}
\label{sec:v-loops}

Слово \code{for} в Verilog означает не «повторять во времени», а «размножить
аппаратуру». Синтезатор раскрывает цикл на этапе сборки: цикл по восьми битам
превращается в восемь одинаковых кусков схемы, которые работают
\emph{одновременно}. Поэтому число повторений должно быть известно заранее
(константа или параметр), а переменная цикла объявляется как \code{integer}
или \code{genvar}.

Есть два вида циклов. \code{for} внутри \code{always} повторяет операторы
(например, суммирует биты). Цикл \code{generate for} повторяет
\emph{объявления}: экземпляры модулей, \code{assign}, \code{always}. Он нужен,
когда схему из $N$ одинаковых ступеней надо описать один раз, а $N$ менять
параметром. \code{function} описывает комбинационную «подпрограмму» без
времени и без памяти: её вызывают в правой части \code{assign}.

\vfile{\codepath/05_verilog/loops.v}{code/05\_verilog/loops.v}

\textbf{Как читается.} В \code{popcount8} переменная \code{n} обнуляется,
затем к ней по очереди прибавляется каждый бит слова: на схеме получится цепочка
сумматоров. В \code{adder\_gen} метка \code{stage} называет каждую ступень цикла
(в иерархии они видны как \code{stage[0]}, \code{stage[1]}\ldots), а провода
переноса \code{c} связывают ступени друг с другом, как в сумматоре из
главы~\ref{ch:arith}. В \code{parity\_fn} функция вычисляет чётность: на RTL это
цепочка XOR.

\begin{diybox}[Цикл и параметр]
\textbf{Задание.} Проверить, что цикл раскрывается в аппаратуру, и менять разрядность параметром.\\
\textbf{В Logra.} Вкладка «Код»: \code{loops.v}; «Синтезировать»; выберите верхним модулем \code{adder\_gen} и измените в его заголовке \code{parameter N = 4} на \code{parameter N = 8}.\\
\textbf{Ожидаемый результат.} Для \code{N = 8} на RTL вдвое больше ступеней; на диаграмме \code{a = 100}, \code{b = 100}, \code{cin = 1} даёт \code{s = 201}.\\
\textbf{Как проверить.} Тестбенч \code{tb\_loops.v} (\code{iverilog -g2012 -o tb tb\_loops.v loops.v}, затем \code{vvp tb}) перебирает все слова и печатает \code{loops: OK}.
\stars{2} Допишите модуль, считающий единицы в 16-разрядном слове.
\end{diybox}

\textbf{Образец ответа.}\par\nopagebreak
\begin{vcode}
module popcount16 (
    input      [15:0] d,
    output reg [4:0]  n         // от 0 до 16: нужно 5 бит
);
    integer i;
    always @(*) begin
        n = 5'd0;
        for (i = 0; i < 16; i = i + 1)
            n = n + d[i];
    end
endmodule
\end{vcode}


\section{Иерархия: модуль внутри модуля}

Большие схемы собирают из маленьких модулей, как устройство из микросхем.
Подключение одного модуля внутри другого называется \term{созданием экземпляра}
(instantiation).

\begin{figure}[htbp]
\centering
\begin{minipage}{0.46\textwidth}
\tcbset{codestyle/.append style={unbreakable}}
\begin{vcode}
module full_adder (
    input  wire a, b, cin,
    output wire s, cout
);
    wire s1, c1, c2;
    half_adder ha1 (.a(a),  .b(b),
                    .s(s1), .c(c1));
    half_adder ha2 (.a(s1), .b(cin),
                    .s(s),  .c(c2));
    assign cout = c1 | c2;
endmodule
\end{vcode}
\end{minipage}\hfill
\begin{minipage}{0.50\textwidth}
\centering
\begin{tikzpicture}[ha/.style={draw=FpgaBlue, fill=FpgaLight, thick, minimum width=13mm, minimum height=13mm, font=\sffamily\footnotesize, align=center}]
  \node[ha] (h1) at (0,0.9) {ha1};
  \node[ha] (h2) at (2.2,-0.6) {ha2};
  \node[draw, circle, thick, minimum size=8mm, font=\scriptsize] (or) at (4.4,0.5) {$\ge1$};
  \draw[wire] ([yshift=3mm]h1.west) -- ++(-0.65,0) node[left] {a};
  \draw[wire] ([yshift=-3mm]h1.west) -- ++(-0.65,0) node[left] {b};
  \draw[wire] ([yshift=-3mm]h2.west) -- (-1.3,-0.9) node[left] {cin};
  \draw[wire] ([yshift=-3mm]h1.east) -- node[above, lbl] {s1} (1.1,0.6) |- ([yshift=3mm]h2.west);
  \draw[wire] ([yshift=3mm]h1.east) -- node[above, lbl] {c1} (3.7,1.2) |- (or.150);
  \draw[wire] ([yshift=3mm]h2.east) -- node[below, lbl] {c2} (3.4,-0.3) |- (or.210);
  \draw[wire] ([yshift=-3mm]h2.east) -- ++(2.35,0) node[right] {s};
  \draw[wire] (or.east) -- ++(0.4,0) node[right] {cout};
  \node[draw=FpgaGray, dashed, fit=(h1)(h2)(or), inner sep=5mm, label={[lbl]above:full\_adder}] {};
\end{tikzpicture}
\end{minipage}
\caption{Полный сумматор из двух полусумматоров}
\end{figure}

Запись \code{.a(s1)} означает: «вход \code{a} экземпляра соединить с проводом
\code{s1}».

\begin{ideabox}
\begin{itemize}
  \item \code{assign} и \code{always @(*)} описывают \textbf{комбинационную}
        логику (элементы И, ИЛИ, мультиплексоры, сумматоры);
  \item \code{always @(posedge clk)} описывает \textbf{триггеры} и всё, что
        перед ними;
  \item все блоки работают \textbf{одновременно}.
\end{itemize}
\end{ideabox}

\begin{diybox}[Полный сумматор из двух модулей]
\textbf{Задание.} Соберите \code{full\_adder} из двух \code{half\_adder} (листинг раздела).\\
\textbf{В Logra.} «Код»: два модуля в одном файле (верхний --- \code{full\_adder}); урок «Модуль внутри модуля»; «Синтезировать»; RTL, Gowin, «Диаграммы».\\
\textbf{Ожидаемый результат.} На RTL видны два вложенных блока \code{ha1} и \code{ha2}; сумма и перенос по диаграмме совпадают с примером «Полный сумматор» из каталога.\\
\textbf{Как проверить.} Перебрав входы, убедитесь, что \code{cout, s} равны двоичной сумме \code{a + b + cin}. \stars{3} \textbf{Этап проекта «Охранная панель»:} вынесите в отдельный модуль логику датчиков (например, «сработал любой из двух»), подключите её в верхний модуль экземпляром. Критерий приёмки: сборка без ошибок, иерархия видна на RTL, на «Диаграммах» все комбинации датчиков проверены.
\end{diybox}

\textbf{Образец ответа к \stars{3}.}\par\nopagebreak
\begin{vcode}
module any_sensor (
    input  a, b,
    output y
);
    assign y = a | b;                   // сработал любой из двух
endmodule

module panel_top2 (
    input  door, window,
    output alarm
);
    any_sensor s1 (.a(door), .b(window), .y(alarm));
endmodule
\end{vcode}


\section{Регистры, сдвиг и делитель}
\label{sec:v-regs}

Все последовательные узлы пишутся по одному шаблону:
\code{always @(posedge clk)}, внутри него сброс, затем условия, затем
присваивание через \code{<=}. Если условие не выполнено, регистр просто хранит
старое значение: в синхронном блоке это нормально и защёлки не создаёт. Ниже
три узла, которые встретятся дальше в книге: регистр с разрешением записи,
сдвиговый регистр и делитель частоты.

\vfile{\codepath/05_verilog/regs.v}{code/05\_verilog/regs.v}

\textbf{Как читается.} \code{reg\_en}~--- регистр из \code{W} триггеров: по
фронту такта он очищается при \code{rst}, записывает \code{d} при \code{en} и
иначе хранит значение (на схеме это триггер с входом разрешения).
\code{shift4} сдвигает содержимое влево на одну позицию и вдвигает справа
новый бит: после четырёх тактов в регистре оказываются четыре последних бита
входа. \code{clk\_div} считает такты до \code{N-1} и на один такт поднимает
\code{tick}: так из частоты 27~МГц получают медленные события (например, 1~Гц
при \code{N = 27\_000\_000}). Параметры \code{W} и \code{N} можно менять при
создании экземпляра, как в тестбенче \code{tb\_regs.v}.

\begin{diybox}[Регистр, сдвиг, делитель]
\textbf{Задание.} Проверить три узла и дописать свой.\\
\textbf{В Logra.} Вкладка «Код»: \code{regs.v}; верхним по очереди выберите
\code{reg\_en}, \code{shift4}, \code{clk\_div}; «Синтезировать»; «Диаграммы»:
число циклов 20, импульс \code{rst} в начале.\\
\textbf{Ожидаемый результат.} \code{reg\_en} меняет \code{q} только при
\code{en = 1}; \code{shift4} после битов 1, 0, 1, 1 показывает \code{1011};
\code{tick} из \code{clk\_div} при \code{N = 4} приходит каждый четвёртый такт.\\
\textbf{Как проверить.} Тестбенч \code{tb\_regs.v} печатает \code{regs: OK}.
\stars{2} Допишите сдвиговый регистр с параллельной загрузкой: при
\code{load = 1} в него записывается слово \code{d}.
\end{diybox}

\textbf{Образец ответа.}\par\nopagebreak
\begin{vcode}
module shift4_load (
    input            clk, rst, load, din,
    input      [3:0] d,
    output reg [3:0] q
);
    always @(posedge clk) begin
        if (rst)       q <= 4'd0;
        else if (load) q <= d;                // загрузка важнее сдвига
        else           q <= {q[2:0], din};    // иначе сдвиг
    end
endmodule
\end{vcode}


\section{Тестбенч: проверка без рук}
\label{sec:v-tb}

До сих пор мы проверяли схемы вручную: подавали входы на вкладке
«Диаграммы» и смотрели выходы. Но руками все комбинации не переберёшь, и
проверку нельзя повторить после каждой правки. Поэтому рядом с модулем пишут
\term{тестбенч} (testbench): отдельный модуль без портов, который сам подаёт
входы, ждёт и сравнивает выходы с ожидаемыми. Он нужен только в симуляции и
на плату не попадает, поэтому в нём разрешено то, что нельзя синтезировать:
\code{initial}, задержки \code{\#}, циклы с переменной границей,
\code{\$display}. Подробнее тестбенчи разобраны в главе~\ref{ch:sim}; здесь
минимум, чтобы читать примеры этой книги.

\vfile{\codepath/05_verilog/tb_half_adder.v}{code/05\_verilog/tb\_half\_adder.v}

\textbf{Как читается.}
\begin{itemize}
  \item \code{`timescale 1ns/1ps} задаёт единицу времени: \code{\#10} значит 10~нс.
  \item Входы проверяемого модуля в тестбенче~--- \code{reg} (их значения
        задаёт тестбенч), выходы~--- \code{wire} (их ведёт модуль).
  \item \code{half\_adder dut (...)} создаёт экземпляр проверяемого модуля
        (\emph{device under test}).
  \item \code{initial begin ... end} выполняется один раз с нулевого
        момента; операторы внутри идут \emph{по порядку}, а \code{\#10}
        останавливает время.
  \item \code{!==} сравнивает и с неизвестными значениями \code{x} и
        \code{z}: обычное \code{!=} при \code{x} давало бы «неизвестно», и
        ошибка могла бы пройти незамеченной.
  \item \code{\$display} печатает строку; \code{\$finish} завершает
        симуляцию. В тестбенчах книги последняя строка всегда
        \code{имя: OK}: по ней видно, что проверка прошла.
\end{itemize}

Генератор такта пишут одной строкой: \code{always \#5 clk = \textasciitilde clk;}~---
период 10~нс. Все тестбенчи этой книги лежат рядом с модулями (файлы
\code{tb\_*.v}) и запускаются так:
\begin{vcode}[text]
iverilog -g2012 -o tb tb_half_adder.v half_adder.v
vvp tb
\end{vcode}

\begin{diybox}[Свой тестбенч]
\textbf{Задание.} Написать тестбенч к модулю и поймать в нём ошибку.\\
\textbf{В Logra.} Вкладка «Код»: \code{half\_adder.v}; рядом положите
\code{tb\_half\_adder.v}. Для проверки в приложении замените его вкладкой
«Диаграммы» $\to$ «Перебрать входы» (она создаёт такой тестбенч сама).\\
\textbf{Ожидаемый результат.} В терминале \code{half\_adder: OK}; на
диаграмме все четыре строки таблицы полусумматора.\\
\textbf{Как проверить.} Намеренно замените в модуле \code{\^{}} на \code{|} и
запустите тестбенч снова: он должен напечатать \code{ОШИБОК} и номер
комбинации, на которой выход неверен.
\stars{2} Напишите тестбенч для \code{full\_adder} из раздела про иерархию.
\end{diybox}

\textbf{Образец ответа.}\par\nopagebreak
\begin{vcode}
`timescale 1ns/1ps
module tb_full_adder;
    reg  a, b, cin;
    wire s, cout;
    integer i, bad = 0;
    full_adder dut (.a(a), .b(b), .cin(cin), .s(s), .cout(cout));
    initial begin
        for (i = 0; i < 8; i = i + 1) begin
            {a, b, cin} = i[2:0];
            #10;
            if ({cout, s} !== a + b + cin) bad = bad + 1;
        end
        if (bad == 0) $display("full_adder: OK");
        else          $display("full_adder: ОШИБОК %0d", bad);
        $finish;
    end
endmodule
\end{vcode}
Для запуска к тестбенчу добавляют файлы \code{full\_adder.v} и \code{half\_adder.v}.


\section{Типичные ошибки на Verilog}
\label{sec:v-pitfalls}

Три ошибки встречаются чаще остальных. Они не приводят к сообщению
«ошибка», схема собирается, но ведёт себя не так, как задумано.

\begin{enumerate}
  \item \textbf{Латч из-за неполного \code{if} или \code{case}.} Если при
        каком-то условии комбинационному блоку не присвоено значение, синтезатор
        добавляет защёлку, которая «помнит» прежнее.
  \item \textbf{Нет \code{default}.} Без него при неописанных значениях выбора
        получается то же самое: латч или лишняя логика.
  \item \textbf{Блокирующее присваивание в тактируемом блоке.} \code{=} внутри
        \code{always @(posedge clk)} делает значения зависимыми от порядка строк,
        и «одновременные» регистры перестают быть одновременными.
\end{enumerate}

\vfile{\codepath/05_verilog/pitfalls.v}{code/05\_verilog/pitfalls.v}

\textbf{Как читается.} В паре \code{latch\_bad} и \code{latch\_ok} различается
одна строка: \code{else q = 1'b0}. В \code{swap\_bad} вторая строка видит уже
новое значение \code{a}, и оба регистра оказываются равными; в \code{swap\_ok}
правые части вычисляются «до фронта» и значения обмениваются. Тестбенч
\code{tb\_pitfalls.v} показывает обе ошибки в симуляции: латч помнит
единицу, хотя \code{en = 0}, а «обмен» даёт одинаковые \code{a} и \code{b}.

\begin{diybox}[Найдите ошибку]
\textbf{Задание.} Найти и исправить три ошибки в коде.\\
\textbf{В Logra.} Вкладка «Код»: \code{pitfalls.v}; верхним выберите
\code{latch\_bad}, «Синтезировать»; вкладка RTL.\\
\textbf{Ожидаемый результат.} В \code{latch\_bad} на RTL виден элемент LATCH;
в \code{latch\_ok} его нет, только мультиплексор.\\
\textbf{Как проверить.} Тестбенч \code{tb\_pitfalls.v} печатает
\code{pitfalls: OK}, потому что \emph{ожидает} неправильное поведение плохих
модулей; когда исправите их, поведение изменится. \stars{2} Найдите
латч в коде ниже и исправьте его.
\end{diybox}

\begin{vcode}
module level_cmp (
    input      [3:0] x,
    output reg       low, mid
);
    always @(*) begin
        if (x < 4'd4)       low = 1'b1;
        else if (x < 4'd8)  mid = 1'b1;
    end
endmodule
\end{vcode}
\textbf{Образец ответа.} Обе переменные должны получать значение в каждой
ветке: сначала \code{low = 0; mid = 0;}, затем \code{if} меняет нужную.
\begin{vcode}
module level_cmp_ok (
    input      [3:0] x,
    output reg       low, mid
);
    always @(*) begin
        low = 1'b0;
        mid = 1'b0;                  // значения по умолчанию: латча нет
        if (x < 4'd4)       low = 1'b1;
        else if (x < 4'd8)  mid = 1'b1;
    end
endmodule
\end{vcode}


\section{Что дальше}

Мы умеем описать на Verilog вентили, мультиплексор и сумматор и проверить их на
вкладках RTL и Диаграммы. Но схемы пока даны готовыми. Как из словесного условия
(«тревога, если дверь открыта и сейчас ночь») получить таблицу истинности и
минимальную схему и не сделать лишних элементов? Этим мы займёмся в
главе~\ref{ch:comb}: на каждом шаге будем строить схему и в Конструкторе, и
текстом на Verilog.
